\section{The proposal and its logical premises}

\subsection{The three relative classes and the two meanings of fibration}

We read the supplied transcriptions of \emph{Les D\'erivateurs} against the supplied edition PDFs wherever those PDFs are present \citep{BISITEGrothendieck1991}. There is no edition PDF for VI-b or XIIc-a in the supplied archive; we additionally opened the editors' public XIIc PDF and checked ms.~207--208, 217--218 and 222--224, including images of the crucial pages. A correction at a location for which only a transcription was checked is justified by typing or by the mathematical argument, rather than claimed as an independent palaeographical verification. All categories are small in the fixed universe, except that a category of models may be large in a second universe; its slice localizations are assumed to exist in that second universe. No local-smallness conclusion follows from two-out-of-three alone.

\begin{claim}{A.1}{Proved}{The historical and adopted localizer axioms are different formulations}
The literal Loc~(1) of XVI ms.~1 consists of isomorphisms, two-out-of-three, and the condition that $gf,fg\in W$ imply $f,g\in W$. Grothendieck calls the first two conditions \emph{saturation faible}, and all three \emph{saturation mod\'er\'ee}. Retract closure is added in Loc~(1 bis); strong saturation is Loc~(1 ter). The literal Loc~(3), XVI ms.~2, is absolute Theorem~A. Relative Theorem~A is Loc~(3 bis), ms.~3.

In this paper, a fundamental localizer means the modern relative package specified in the conventions: weak saturation, the terminal-object axiom, and the assertion that a functor $u:A\to B$ over $C$ lies in $W$ whenever $u/c:A/c\to B/c$ lies in $W$ for every $c\in C$. Consequently no argument using relative Theorem~A is an argument from the literal absolute Loc~(3) alone.
\end{claim}
\begin{proof}
These are the displayed clauses and terminology of XVI ms.~1--3, checked in both the transcription and the edition PDF. In the relative axiom take $B=C$ and its structure map to be the identity. Since $C/c$ has terminal object $1_c$, the terminal-object axiom and two-out-of-three identify its premise with $A/c\to e\in W$. Thus relative Theorem~A implies the literal absolute form. The reverse implication is not a formal substitution.
\end{proof}

\begin{claim}{A.2}{Proved}{Intrinsic fibrations and their unconditional closure properties}
Let $\mathcal M$ have pullbacks and let $W$ be invariant under isomorphism of arrows. For $S\in\mathcal M$, let $W_S$ be the inverse image of $W$ under the source functor $\mathcal M/S\to\mathcal M$. Define
\[
\operatorname{Fib}_W=\{p:E\to B\mid p^*: \mathcal M/B\to\mathcal M/E\text{ sends }W_B\text{ into }W_E\}.
\]
Thus every test is an arrow $v:T'\to T$ over $B$, with arbitrary structure map $T\to B$, and the tested arrow is $E\times_B T'\to E\times_B T$. This class contains isomorphisms and is closed under composition and arbitrary pullback. These assertions require no two-out-of-three axiom. If $W$ is closed under retracts, then so is $\operatorname{Fib}_W$.
\end{claim}
\begin{proof}
For $E\xrightarrow p B\xrightarrow q C$, the pullback of an arrow over $C$ along $qp$ is canonically isomorphic to its successive pullbacks along $q$ and $p$. Both preservation hypotheses apply in that order. For a pullback $p':E\times_B B'\to B'$, forget a test arrow over $B'$ to an arrow over $B$ and use the canonical associativity isomorphism of pullbacks. An isomorphism induces an equivalence of slices whose underlying arrows are isomorphic to the originals. Finally, if $p$ is a retract of $q$ in the arrow category, the pullback of any fixed test arrow along $p$ is a retract of the corresponding pullback along $q$, after pulling the latter over the retraction maps of bases. The two maps exhibiting this retract are induced by the maps in the arrow-category retract diagram. Retract closure of $W$ now applies. The additional hypothesis cannot be suppressed merely because XII ms.~208 lists direct factors in a parenthesis: that paragraph has explicitly assumed only isomorphisms and two-out-of-three.
\end{proof}

\begin{claim}{A.3}{Proved}{The coarse, universal, and local relative classes}
Suppose $W$ contains isomorphisms and satisfies two-out-of-three. Let $\mathcal L$ contain isomorphisms and be closed under composition and pullback. For an arrow $u:X\to Y$ in $\mathcal M/S$, put
\[
u\in W_S^{\mathcal L}\quad\Longleftrightarrow\quad
\bigl(X\times_S T\longrightarrow Y\times_S T\bigr)\in W
\quad\hbox{for every }(T\to S)\in\mathcal L.
\]
Let $W_S^u$ be the same class with all $T\to S$ allowed. Then
\[
W_S^u\subseteq W_S^{\mathcal L}\subseteq W_S.
\]
All three classes contain isomorphisms and satisfy two-out-of-three. Their retract closure follows if $W$ has it. These are precisely the relative classes in XII ms.~210 and 220. They must not be confused with the class of universal $W$-equivalences in $\mathcal M$, where the base is the codomain of the arrow being tested.
\end{claim}
\begin{proof}
Each test is a functor on the slice. The asserted class is the intersection of the inverse images of $W$ under those functors; identities, two-out-of-three, and retracts are preserved by this operation. The identity $S\to S$ is a permitted local test, which proves the second inclusion; the first follows by reducing the family of tests. In particular no descent argument can replace the two-out-of-three premise. A class defined only by full faithfulness of a restriction functor may fail it, as the set-coefficient example $e\to BC_2\to e$ does: restriction along $BC_2\to e$ and along their composite is fully faithful, whereas forgetting the $C_2$-action along $e\to BC_2$ is not full.
\end{proof}

\begin{claim}{A.4}{Proved}{What raw reindexing and postcomposition actually preserve}
Under the hypotheses of A.3, a map $f:S'\to S$ in $\mathcal L$ induces
\[
f^*:(\mathcal M/S,W_S^{\mathcal L})\longrightarrow
(\mathcal M/S',W_{S'}^{\mathcal L}).
\]
For every $f:S'\to S$, without assuming $f\in\mathcal L$, postcomposition induces
\[
f_\#:(\mathcal M/S',W_{S'}^{\mathcal L})\longrightarrow
(\mathcal M/S,W_S^{\mathcal L}).
\]
If $f^*$ also preserves the displayed classes, the ordinary adjunction $f_\#\dashv f^*$ descends to their localizations. For arbitrary $f$, this calculation alone constructs the descended postcomposition functor; it does not construct its right adjoint.
\end{claim}
\begin{proof}
For the first assertion, test $f^*u$ along $g:T\to S'$ in $\mathcal L$. Its underlying arrow is the pullback of $u$ along $fg\in\mathcal L$. For the second, test $f_\#u$ along $h:T\to S$ in $\mathcal L$. The projection $S'\times_S T\to S'$ lies in $\mathcal L$, and the tested underlying arrow is exactly the pullback of $u$ along this projection. For descent of the adjunction, apply the localization functors to the original unit and counit. Their naturality extends across inverted arrows, and their two triangle identities remain identities. Thus they exhibit the descended adjunction. This also proves the general adjunction lemma at XII ms.~248 and settles the formal postcomposition observation at XII ms.~226--227 with its exact scope.
\end{proof}

\begin{claim}{A.5}{Proved}{The universal saturation of the local test class}
Under A.3, define $\widehat{\mathcal L}$ by the following condition on $f:S'\to S$: for every $T\to S$, pullback along $f_T:T\times_S S'\to T$ sends $W_T^{\mathcal L}$ into $W_{T\times_S S'}^{\mathcal L}$. Then $\widehat{\mathcal L}$ contains $\mathcal L$, contains isomorphisms, and is closed under composition and pullback. Moreover
\[
W_S^{\widehat{\mathcal L}}=W_S^{\mathcal L}
\quad\hbox{for every }S,\qquad
\widehat{\widehat{\mathcal L}}=\widehat{\mathcal L}.
\]
It is the largest composition- and pullback-closed class containing isomorphisms that defines this same family of fine relative classes. This is the corrected, second definition at XII ms.~222; the first definition on that page omits the universal quantifier over $T\to S$.
\end{claim}
\begin{proof}
A.4 proves $\mathcal L\subseteq\widehat{\mathcal L}$, also after any preliminary pullback. Iterated pullback proves closure under composition; absorbing a preliminary pullback into $T\to S$ proves pullback closure. If $u\in W_S^{\mathcal L}$ and $f\in\widehat{\mathcal L}$ has target $S$, its pullback lies in $W_{S'}^{\mathcal L}$, hence in $W_{S'}$, which proves $W_S^{\mathcal L}\subseteq W_S^{\widehat{\mathcal L}}$. The reverse inclusion follows from $\mathcal L\subseteq\widehat{\mathcal L}$. The defining condition for a hat depends only on this family, proving idempotence. If $\mathcal K$ has the stated closure properties and defines the same family, apply A.4 to $\mathcal K$ after any preliminary pullback: every member of $\mathcal K$ then satisfies the defining condition of $\widehat{\mathcal L}$.
\end{proof}

\subsection{Transporters, stabilizers, and the cancellation in Chapter II}

\begin{claim}{A.6}{Proved}{The corrected transporter Galois connection}
Let $W$ contain isomorphisms and be closed under composition in a category $\mathcal M$ with pullbacks. Let $A\subseteq W$ contain isomorphisms and be composition-closed, and let $L$ contain isomorphisms and be pullback-closed. Define
\[
\begin{split}
T_W(A)&=\{p:E\to B\mid p^*\text{ sends every arrow of }A
\text{ in }\mathcal M/B\text{ into }W\},\\
K_W(L)&=\{u:X\to Y\in W\mid p^*u\in W
\text{ for every }(p:Y'\to Y)\in L\}.
\end{split}
\]
Then $T_W(A)$ contains isomorphisms and is pullback-closed, $K_W(L)$ contains isomorphisms and is composition-closed, and
\[
A\subseteq K_W(L)\quad\Longleftrightarrow\quad L\subseteq T_W(A).
\]
Both operations reverse inclusions; their composites are extensive, monotone, and idempotent. The closed pairs are exchanged by inverse order-reversing bijections. Their extreme pairs are
\[
(W^u,\operatorname{Mor}\mathcal M),\qquad
(W,\operatorname{Fib}_W),
\]
where $W^u$ means universal equivalences with respect to the codomain of each arrow.
\end{claim}
\begin{proof}
Pullback stability of $T_W(A)$ uses the quantifier over arrows in a slice: a test over $B'$ may be regarded as a test over $B$. To show composition stability of $K_W(L)$, pull back a composite $vu$ along $p\in L$. The resulting composite consists of the pullback of $v$ along $p$ and the pullback of $u$ along the pullback of $p$ to the source of $v$. Both belong to $W$.

If $L\subseteq T_W(A)$, apply its defining preservation condition to an arrow of $A$ viewed over its own codomain. Conversely, suppose $A\subseteq K_W(L)$, and let $u:X\to Y$ in $A$ have a further structure map $Y\to B$. For $p:E\to B$ in $L$, the pullback $E\times_B Y\to Y$ lies in $L$; the kernel test for $u$ therefore gives $p^*u\in W$. This proves the adjunction. Antitonicity is immediate from enlarging the test families. The familiar threefold identities follow directly: extensivity gives $T_WK_WT_W(A)\supseteq T_W(A)$, while applying antitonicity to $A\subseteq K_WT_W(A)$ gives the reverse inclusion. The kernel identity is identical with the roles reversed. These identities imply idempotence. Finally $T_W(\operatorname{Iso})$ is all arrows, $K_W(\operatorname{Mor}\mathcal M)=W^u$, $K_W(\operatorname{Iso})=W$, and $T_W(W)=\operatorname{Fib}_W$. These arguments verify the corrected domains of II~(6.7), ms.~86; composition must not be imposed on the right-hand domain in advance.
\end{proof}

\begin{claim}{A.7}{Proved}{The exact composition criterion, II Proposition 6.1}
For a closed pair $A=K_W(L)$, $L=T_W(A)$ of A.6, the following conditions are equivalent: $L$ is composition-closed; pullback along every member of $L$ preserves $A$; and $L=\operatorname{Stab}(A)$, where $\operatorname{Stab}(A)=T_A(A)$. Every stabilizer is composition- and pullback-closed. In particular, a transporter into the larger class $W$ is not generally a stabilizer.
\end{claim}
\begin{proof}
If $L$ composes, pull an arrow $u\in K_W(L)$ back along $p\in L$, and test the result along $q\in L$. The resulting arrow is a test of $u$ along the composite of $q$ with the appropriate pullback of $p$, which is again in $L$. Thus $p^*u\in A$. Conversely, if $L$ preserves $A$, pulling an $A$-map successively along two members of $L$ again gives an $A$-map and hence a $W$-map. The composite therefore lies in $T_W(A)=L$. Preservation gives $L\subseteq\operatorname{Stab}(A)$; the other inclusion follows from $A\subseteq W$. Finally stabilization under pullback and composition is the same direct argument as A.2 with $A$ in place of $W$.
\end{proof}

\begin{claim}{A.8}{Refuted}{Transporters need not compose under the hypotheses of II section 6}
The claim in the first version of II~(6.2)--(6.5) that $T_W(A)$ automatically belongs to a right-hand domain requiring composition closure is false. A counterexample is the chain category $0<1<2<3$, with $W$ consisting of the identities and $0\to2,0\to3$, and $A$ consisting of the identities and $0\to3$. Here
\[
T_W(A)=\operatorname{Mor}([3])\setminus\{1\to3\},\qquad
K_WT_W(A)=A.
\]
Consequently $1\to2$ and $2\to3$ lie in the transporter while their composite does not.
\end{claim}
\begin{proof}
There is exactly one arrow $i\to j$ when $i\leq j$, and composition is forced by the order. Pullbacks are minima. Both displayed classes contain all isomorphisms and are closed under composition. The only nonidentity test in $A$ is $0\to3$, and it can be regarded as an arrow over a base $b$ only for $b=3$. Pulling it back along $i\to3$ gives $0\to i$, which lies in $W$ exactly for $i=0,2,3$. This proves the transporter formula including its slice quantifiers. Of the two nonidentity arrows in $W$, the arrow $0\to2$ fails the kernel test along $1\to2$, since its pullback is $0\to1\notin W$. The other survives by the Galois relation. The finite checker \texttt{code/a\_finite\_audit.py} verifies the category, all pullbacks, both computed classes, and the failed composite. We do not claim that this $W$ satisfies two-out-of-three; that axiom is not among the hypotheses of II section 6.
\end{proof}

\begin{claim}{A.9}{Proved}{The meaning of smooth and proper in the manuscript}
For a fundamental localizer on $\mathrm{Cat}$, let $\operatorname{As}_W$ consist of $u:A\to B$ for which $A/b\to e$ lies in $W$ for every $b$, and let $\operatorname{Coas}_W$ consist of $u$ for which $b\backslash A\to e$ lies in $W$ for every $b$. The definitions used at II ms.~66 and explicitly reaffirmed at ms.~89 are
\[
\operatorname{Liss}_W=\operatorname{Stab}(\operatorname{As}_W),
\qquad
\operatorname{Prop}_W=\operatorname{Stab}(\operatorname{Coas}_W).
\]
The two mixed transporters contain both the corresponding stabilizer and $\operatorname{Fib}_W$. Neither is the definition of smoothness or properness. The proposed inclusion
$\operatorname{Liss}_W\cap\operatorname{Prop}_W\subseteq\operatorname{Fib}_W$
at XII ms.~225--226 is explicitly qualified by ``peut-\^etre'' and is an extra proposed axiom, not a consequence asserted to have been proved there.
\end{claim}
\begin{proof}
The two definitions are written in II~(4.10) and in the paragraph introducing stabilizers after Proposition~6.1. Absolute Theorem~A gives $\operatorname{As}_W\subseteq W$; opposition gives the analogous containment for $\operatorname{Coas}_W$. Therefore an $\operatorname{As}_W$-preserving functor sends such maps into $W$, and a $W$-preserving functor does so as well. These are the two claimed containments; the coaspherical argument is identical. XII ms.~225 introduces two arbitrary classes, asks whether to assume these closures and the parenthesized inclusion, and then separately proposes the fine-localization axioms. This wording does not identify the inclusion with an earlier theorem. The finite lifting criteria and the actual validity range of the inclusion are treated in Part B.
\end{proof}

\subsection{The fibre axioms and the later implication scheme}

\begin{claim}{A.10}{Proved}{The three fibre axioms, with their quantifiers}
For every $S\in\mathcal M$ let $\mathcal F_S$ be the full subcategory of $\mathcal M/S$ on the intrinsic fibrations. The three assertions at XII ms.~215--218 are
\[
\begin{aligned}
\mathrm{(a)}&\quad J_S:\mathcal F_S[(W_S^u)^{-1}]
 \longrightarrow(\mathcal M/S)[W_S^{-1}]\text{ is an equivalence};\\
\mathrm{(b)}&\quad W_S^u=W_S\text{ on }\mathcal F_S;\\
\mathrm{(c)}&\quad\mathcal F_S[W_S^{-1}]\longrightarrow
 \mathcal H[W_S^{-1}]\longrightarrow(\mathcal M/S)[W_S^{-1}]
 \text{ are equivalences}
\end{aligned}
\]
for every full $\mathcal F_S\subseteq\mathcal H\subseteq\mathcal M/S$.
Axiom (a) implies equality of the strong saturations of $W_S^u$ and $W_S$ on $\mathcal F_S$. It does not by that argument imply equality of the original classes. Axioms (b) and (c) imply (a).
\end{claim}
\begin{proof}
The displayed formulations are exactly the three source formulas, with all localizations and restrictions typed. If $u\in W_S$ has source and target in $\mathcal F_S$, then $J_S$ sends its image to an isomorphism. An equivalence of categories reflects isomorphisms, so $u$ is inverted by localization at $W_S^u$. The reverse containment follows from A.3. This is a statement about strong saturation, not about membership in $W_S^u$. For the last assertion apply (c) to either endpoint and use (b) to identify the source localization with the source of $J_S$.
\end{proof}

\begin{claim}{A.11}{Proved}{Enough intrinsic fibrations implies the exact fibrewise invariance}
Suppose $\mathcal M$ has pullbacks, $W$ contains isomorphisms and satisfies two-out-of-three, and every map factors as a $W$-map followed by a member of $\operatorname{Fib}_W$. Then axiom (b) of A.10 holds, and
\[
\operatorname{Fib}_W\cap W=W^u.
\]
No terminal object, accessibility, functorial factorization, or strong saturation is required for these conclusions.
\end{claim}
\begin{proof}
Let $u:X\to Y$ be a $W$-map over $S$ with $X\to S$ and $Y\to S$ intrinsic fibrations. Factor an arbitrary base change as $T\xrightarrow{i}T_0\xrightarrow{p}S$, with $i\in W$ and $p\in\operatorname{Fib}_W$. The pullback $u_0:X\times_S T_0\to Y\times_S T_0$ lies in $W$ because $p$ is intrinsic. The two vertical comparison maps from pullback over $T$ to pullback over $T_0$ lie in $W$, because the pullbacks of $X\to S$ and $Y\to S$ are intrinsic and preserve $i$. The resulting commutative square and two-out-of-three imply $u_T\in W$. This proves (b).

For $q:X\to S$ in $\operatorname{Fib}_W\cap W$, apply (b) to the arrow from $q$ to $1_S$, both fibrant objects of the slice; thus $q\in W^u$. Conversely suppose $q\in W^u$ and let $v:T'\to T$ be a $W$-map over $S$. The vertical maps in
\[
\begin{matrix}
X\times_S T'&\longrightarrow&X\times_S T\\
\downarrow&&\downarrow\\
T'&\xrightarrow{v}&T
\end{matrix}
\]
are in $W$ by universal equivalence. Two-out-of-three implies that the upper arrow lies in $W$. Hence $q\in\operatorname{Fib}_W$, and the identity base change gives $q\in W$. This verifies the formal fibration half of XIII Proposition~4, ms.~10--11, and the argument used in XVI~4.4.1.
\end{proof}

\begin{claim}{A.12}{Proved}{A functorial replacement proves every intermediate-subcategory comparison}
Let $\mathcal A$ be a category with a class $W$ containing isomorphisms and satisfying two-out-of-three, let $\mathcal F\subseteq\mathcal A$ be full, and suppose there is $R:\mathcal A\to\mathcal F$ with a natural transformation $\eta:1_{\mathcal A}\to iR$ whose components are in $W$. For every full $\mathcal F\subseteq\mathcal H\subseteq\mathcal A$, both inclusions induce equivalences
\[
\mathcal F[W^{-1}]\longrightarrow\mathcal H[W^{-1}]
\longrightarrow\mathcal A[W^{-1}].
\]
In particular, such slice replacements into $\operatorname{Fib}_W S$, together with enough intrinsic fibrations, verify XII (a), (b), and (c).
\end{claim}
\begin{proof}
For $u:X\to Y$ in $W$, the equation $R(u)\eta_X=\eta_Yu$ implies $R(u)\in W$ by two-out-of-three. Thus $R$ is a relative functor. Its restriction to any $\mathcal H$ has image in $\mathcal F\subseteq\mathcal H$, and each component of $\eta$ at an object of $\mathcal H$ is a morphism of $\mathcal H$ because that subcategory is full. After localization, these components give a natural isomorphism between the identity and inclusion followed by $R$. On $\mathcal F$ they give the other natural isomorphism. The same proof works with $\mathcal A$ in place of $\mathcal H$. Finally A.11 identifies the two classes on fibrant objects, so the claimed source in (a) is the localization used here.
\end{proof}

\begin{claim}{A.13}{Refuted}{The implication asked in audit row XIIc B062}
The axioms (a) and (b) of A.10, even imposed at every base, do not imply (c). There is a counterexample with six objects, all finite limits and colimits, and strongly saturated $W$. It also refutes the analogous assertion for enhanced localizations.
\end{claim}
\begin{proof}
The complete construction and verification are E.12. Its poset has Hasse covers $0<1$, $0<2$, $1<3$, $2<3$, $1<4$, $3<5$, $4<5$, and $W$ consists of identities together with $1\to3$, $1\to4$, $2\to3$. At base $5$, the intrinsic fibrations form the chain $\{0<3<5\}$. Their inclusion into the whole slice is a localization equivalence, certified by explicit relative inverse functors and weak natural transformations. The full intermediate category $\{0,3,4,5\}$ is a diamond with only identity weak arrows, so its localization cannot be equivalent to that chain under the inclusion. E.12 checks (a) and (b) at all other bases and proves strong saturation and the enhanced version. The independent finite checker is \texttt{code/a\_b062\_verify.py}.
\end{proof}
The failure is the missing intermediate-subcategory control: the weak natural zigzag proving (a) passes through an object omitted from the chosen full intermediate subcategory. A.12 states sufficient comparison data that prevent this failure. The example has a terminal object but does not require, or satisfy, the extra condition that every terminal map be an intrinsic fibration; that condition is not part of the setup at XII ms.~207 or the question at ms.~218.

\begin{claim}{A.14}{Proved}{The two fine axiom systems in XII ms.225--227}
Suppose classes $L,P\subseteq\operatorname{Mor}\mathcal M$ contain isomorphisms and are closed under composition and pullback. Define $W_S^g=W_S^L$ and $W_S^d=W_S^P$. The clause called (a) on ms.~226, asserting preservation of $W^g$ by $L$-base change and of $W^d$ by $P$-base change, follows formally from the closure hypotheses. The clause called (b) on that page is the additional pair of assertions
\[
\begin{aligned}
(L/S)[(W_S^u)^{-1}]&\simeq(\mathcal M/S)[(W_S^d)^{-1}],
 &W_S^u=W_S^d\text{ on }L/S,\\
(P/S)[(W_S^u)^{-1}]&\simeq(\mathcal M/S)[(W_S^g)^{-1}],
 &W_S^u=W_S^g\text{ on }P/S.
\end{aligned}
\]
Here $L/S$ and $P/S$ denote full subcategories on the corresponding structural maps. Under these additional equivalences, each fine theory is a contravariant pseudofunctor on all of $\mathcal M$. This conclusion does not require $L\cap P\subseteq\operatorname{Fib}_W$.
\end{claim}
\begin{proof}
The first assertion is A.4. Pullback preserves the objects of $L/S$ and $P/S$ by the stipulated closure, and preserves $W_S^u$ by its universal quantifier. These relative slice categories therefore form a pseudofunctor, with the canonical associativity and unit isomorphisms of pullback; those isomorphisms satisfy their usual coherence identities before localization and hence afterwards. Transporting this pseudofunctor across the displayed equivalences gives the asserted fine pseudofunctors. Neither this construction nor its coherence uses the intersection condition. The source calls a further pair of adjunction properties clause (c), refers to VI ms.~108--113, and asks whether its hypotheses can be simplified. This (c) is a different assertion from the intermediate-subcategory (c) at ms.~218.
\end{proof}

\begin{claim}{A.15}{Proved}{The four fibration classes and Loc (7), with reference localizer}
For a fundamental localizer $W$, the four classes in XVI section~4 are distinct definitions. Besides intrinsic $\operatorname{Fib}_W$, one has Grothendieck fibrations; $\operatorname{Fib}^{\mathrm{weak}}_W$, tested by all arrows $T'\to T$ over the base with $T'$ and $T$ contractible in the category-homotopy sense; and $\operatorname{Parf}_W=\operatorname{Liss}_W\cap\operatorname{Prop}_W$. Here category-homotopy contractibility means that the identity functor is connected to a constant functor by a finite zigzag of natural transformations. It implies weak contractibility of the nerve, but replacing the test family by all weakly contractible nerves is a separate assertion, not the literal definition. Let $\mathcal G_W$ denote Grothendieck fibrations all of whose transport functors lie in $W$. For any arrow class $\mathcal F$, define
\[
\begin{aligned}
\mathrm A(\mathcal F)&:\quad\mathcal F\subseteq\operatorname{Fib}_W,\\
\mathrm B(\mathcal F)&:\quad W_S\cap\operatorname{Mor}(\mathcal F/S)
 \subseteq W_S^u\quad\hbox{for every }S,\\
\mathrm C(\mathcal F)&:\quad\mathcal F\cap W\subseteq W^u.
\end{aligned}
\]
Loc~(7a), (7b), and (7c) are respectively $\mathrm A(\mathcal G_W)$, $\mathrm B(\mathcal G_W)$, and $\mathrm C(\mathcal G_W)$. For a reference localizer $V\subseteq W$, the conditions $\mathrm A_V,\mathrm B_V,\mathrm C_V$ replace $\mathcal G_W$ by $\mathcal G_V$. The condition $\mathrm D_V$ is $\mathrm A(\operatorname{Parf}_V)$, or its equivalent B and C forms once the hypotheses of A.17 are verified. Loc~(7F) is the factorization condition of A.11, not the inclusion of any one of these reference classes in $\operatorname{Fib}_W$.
\end{claim}
\begin{proof}
These are the formulas of XVI~4.1--4.4.1, ms.~27--41. The weak condition tests arbitrary maps between category-contractible bases, not every $W$-map between arbitrary bases. A natural transformation realizes to a homotopy, so a zigzag from the identity to a constant does make the nerve weakly contractible. No converse is used. The class $\mathcal G_W$ adds both the cartesian lifting property and the assertion that each induced transport lies in $W$; neither assertion is part of the definition of intrinsic fibrations. The printed use of ``cofibrations'' in the (7F) paragraph is inconsistent with the factorization and with the displayed class $\operatorname{Fib}_W$ and must read ``fibrations''; the edition PDF contains the same word. Thus the separation is mathematical, not merely typographical.
\end{proof}

\begin{claim}{A.16}{Proved}{The implications A to B to C and replacement transfer}
Let $W$ and $\mathcal M$ satisfy A.11's initial assumptions. If every map factors as a $W$-map followed by a member of $\mathcal F$, then $\mathrm A(\mathcal F)$ implies $\mathrm B(\mathcal F)$. If identities lie in $\mathcal F$, then $\mathrm B(\mathcal F)$ implies $\mathrm C(\mathcal F)$.

Suppose that each $p:X\to S$ in $\mathcal F$ admits a comparison $j_p:X\to R_p$ over $S$ in $W_S^u$, with $R_p\to S$ in $\mathcal G$. Then $\mathrm A(\mathcal G)\Rightarrow\mathrm A(\mathcal F)$ and $\mathrm C(\mathcal G)\Rightarrow\mathrm C(\mathcal F)$. If for every map $h:X\to Y$ over $S$ between $\mathcal F$-objects the chosen comparisons extend to a commutative square $R(h)j_X=j_Yh$, then also $\mathrm B(\mathcal G)\Rightarrow\mathrm B(\mathcal F)$. Functorial replacements suffice; the proof only needs these individual commuting squares.
\end{claim}
\begin{proof}
The first two assertions follow from A.11 and from taking the target object in the slice to be $1_S$. For A-transfer, compare the pullbacks of a $W$-arrow $T'\to T$ over $S$ along $X$ and along $R_p$. Both vertical comparison maps are in $W$ by the universal hypothesis on $j_p$, and the lower horizontal map is in $W$ by $\mathrm A(\mathcal G)$; two-out-of-three gives the upper one. For C-transfer, if $p\in W$, then $R_p\to S$ lies in $W$ by two-out-of-three. It and $j_p$ remain in $W$ after any base change, so their composite does. For B-transfer, if $h\in W$, the commuting square and the two comparisons show $R(h)\in W$. Applying $\mathrm B(\mathcal G)$ makes $R(h)$ universally in $W$, and the same square after arbitrary pullback proves the assertion for $h$. These are the exact content of XVI Lemmas~4.4.1 and~4.4.9 and Corollary~4.4.10, with the compatibility premise stated rather than suppressed.
\end{proof}

\begin{claim}{A.17}{Proved}{The full Brown implication with all its premises}
Suppose $\mathcal M$ has a terminal object and pullbacks, and $W$ contains isomorphisms and satisfies two-out-of-three. Let $\mathcal F$ contain isomorphisms and all terminal maps, be closed under composition and pullback, and admit a factorization of every arrow as a $W$-map followed by an $\mathcal F$-map. Then
\[
\mathrm A(\mathcal F)\quad\Longleftrightarrow\quad
\mathrm B(\mathcal F)\quad\Longleftrightarrow\quad
\mathrm C(\mathcal F).
\]
This is the precise abstract assertion of XVI Lemma~4.4.3. The terminal-map premise is part of the statement, not a consequence of an arbitrary intrinsic-fibration factorization.
\end{claim}
\begin{proof}
Only $\mathrm C\Rightarrow\mathrm A$ remains after A.16. Assume C. Then pullbacks of $\mathcal F\cap W$ are again in $\mathcal F\cap W$. First observe the elementary functor lemma. If a functor $H$ from a category satisfying the stated factorization and terminal-map conditions sends $\mathcal F\cap W$ into a two-out-of-three class $V$, it sends $W$ into $V$. For $u:X\to Y$ in $W$, factor $(1_X,u):X\to X\times Y$ as $w:X\to Z$ in $W$ followed by $q:Z\to X\times Y$ in $\mathcal F$. Both projections $q_X:Z\to X$ and $q_Y:Z\to Y$ lie in $\mathcal F$, since projections are pullbacks of terminal maps. The equations $q_Xw=1_X$ and $q_Yw=u$ show that both lie in $W$. Hence $H(q_X),H(q_Y)\in V$, and the same equations imply $H(w),H(u)\in V$.

For any fixed base $A$, the full category of $\mathcal F$-objects over $A$ satisfies these hypotheses: products are fibre products over $A$, and an absolute factorization is a relative one because $\mathcal F$ composes. Pullback along any $A'\to A$ sends its trivial $\mathcal F$-maps into $W$ by C. The functor lemma therefore proves preservation of every $W$-map between these fibrant objects.

Now let $p:E\to A$ lie in $\mathcal F$ and $u:A'\to A$ in $W$. The factorization just constructed for $u$ writes it as $u=bw$, where $a:Z\to A'$ and $b:Z\to A$ lie in $\mathcal F\cap W$, and $aw=1_{A'}$. Pullback of $b$ is in $W$ by C, so it suffices to handle a base map $w:A'\to A$ admitting a retraction $r:A\to A'$ in $\mathcal F\cap W$.

In that case put $E_1=A\times_{A'}E$, using $r:A\to A'$ and $rp:E\to A'$. Let $l:E_1\to A$ and $k:E_1\to E$ be the projections. The map $l$ lies in $\mathcal F$, and $k$ lies in $\mathcal F\cap W$. There are maps
\[
f=(p,1_E):E\to E_1,
\qquad g=(wrp,1_E):E\to E_1.
\]
They are defined because $rw=1_{A'}$; both retract $k$, so both lie in $W$. The map $f$ is a map over $A$ between the $\mathcal F$-objects $p$ and $l$. Hence $w^*f\in W$ by the relative functor lemma. The canonical identification $A'\times_A E_1\cong E$ identifies its comparison into $E_1$ with $g$. If $q:A'\times_AE\to E$ is the desired pullback of $w$, the pullback square gives $fq=g(w^*f)$. The right side and $f$ lie in $W$, and hence $q\in W$. Applying the argument after a preliminary pullback of $p$ proves the full slice preservation condition $p\in\operatorname{Fib}_W$. This is the mapping-path proof of Brown's lemma, with its relative step made explicit; compare Brown's section~4, Lemmas~1--2 \citep{BISITEBrown1973}.
\end{proof}

\begin{claim}{A.18}{Refuted}{The weak constant-coefficient notion does not compose}
The universal weak rational cohomological fibrations of I ms.~117--118 need not be closed under composition. The failure occurs for finite categories. Consequently that notion cannot be silently substituted for the intrinsic class in A.2 or for the class in Loc~(7F).
\end{claim}
\begin{proof}
Let $G=C_2=\{0,1\}$ under addition. Define $B$ to have objects $a,b$, identity only at $a$, automorphisms $G$ at $b$, arrows $h_g:a\to b$ for $g\in G$, and no arrows $b\to a$. Composition of automorphisms adds their labels and $k_t h_g=h_{t+g}$. Let $q:B\to[1]$ send $a$ to $0$ and $b$ to $1$. It is the Grothendieck opfibration obtained from the functor $e\to BG$. Its fibres are $e$ and $BG$.

For presheaf coefficients, right Kan extension along an opfibration is computed on its strict fibres: the inclusion of the fibre in the over-comma category has a left adjoint given by cocartesian transport. This remains true after any base change. For every complex $M$ of rational vector spaces, the constant-coefficient cohomology of $BG$ is canonically $M$, since the trivial $\mathbb Q[G]$-module $\mathbb Q$ is projective, with the augmentation split by $(1+k_1)/2$. Thus $q_*$ of a constant coefficient is locally constant with the strict-fibre formula, universally under base change.

Define a functor $L:B\to\mathrm{Set}$ by $L(a)=L(b)=G$ and by sending a label $g$ to translation $z\mapsto g+z$. Its category of elements $E$ has objects $(a,z),(b,z)$ for $z\in G$, and arrows $(c,z)\to(d,z')$ given by arrows of $B$ whose assigned translation sends $z$ to $z'$. Composition is the composition in $B$. The projection $p:E\to B$ is a discrete fibration and a discrete opfibration, since each translation is bijective. Right direct images of constants are products over its discrete fibres, and every transport permutes the factors, universally. Thus $p$ has the same weak property.

The composite $qp$ has discrete two-object fibre over $0$ and the contractible regular action groupoid $G\ltimes G$ over $1$. The transport on constant rational cohomology is the diagonal $\mathbb Q\to\mathbb Q^2$. It is not invertible. Hence the direct image under $qp$ of the constant coefficient is not locally constant. The checker lists all arrows, verifies identities and all associative triples, verifies the covering lifts for $p$ and cocartesian lifts for $q$, and computes the two fibre component counts. This independently verifies the finite category used in GALOIS D.3; the averaging argument supplies the unbounded-complex assertion that a component count alone would not establish.
\end{proof}

\begin{claim}{A.19}{Proved}{The local-coefficient repair has an asymmetric conclusion}
Let $\mathbb D$ be a derivator in presheaf variance. Suppose that for $p:E\to B$, after every strict base change, the right direct image of every constant coefficient is locally constant and its value at each base object is canonically the right Kan extension from the strict fibre. If $v:T'\to T$ over $B$ is fully faithful on locally constant $\mathbb D$-coefficients, then the pullback $w:E\times_BT'\to E\times_BT$ is a constant-coefficient $\mathbb D$-equivalence. This argument does not assert full faithfulness or an equivalence on all locally constant coefficients upstairs.

If the same direct-image and base-change hypothesis is required for every locally constant coefficient on every base-changed total category, then the resulting cohomological-fibration class is closed under composition and pullback. The corresponding statements for left Kan extensions and homology follow by duality.
\end{claim}
\begin{proof}
Write $p_T:E_T\to T$ and $p_{T'}:E_{T'}\to T'$ for the projections. For a constant coefficient $M$, put $K=(p_T)_*M_{E_T}$. The canonical mate $v^*K\to(p_{T'})_*M_{E_{T'}}$ is an isomorphism: at each $t'$, the assumed fibre formulas identify it with the identity comparison on the same strict fibre. The object $K$ is locally constant. Full faithfulness of $v^*$ on locally constant objects, applied with source the constant object $N_T$ for every $N\in\mathbb D(e)$, gives
\[
\operatorname{Hom}_{\mathbb D(T)}(N_T,K)
 \cong\operatorname{Hom}_{\mathbb D(T')}(N_{T'},v^*K).
\]
Adjunction to the point and Yoneda identify this with the canonical cohomology comparison; composition of right Kan extensions identifies that comparison with the one induced by $w$ for coefficient $M$. Only constant $M$ on $E_T$ was available, so this proves exactly the stated conclusion.

Under the stronger hypothesis, start with an arbitrary local system $L$ on $E_T$. The same calculation applies because $(p_T)_*L$ is now known to be locally constant. If $E\xrightarrow p B\xrightarrow q S$ are two maps satisfying the stronger hypothesis, then $p_*L$ and $q_*p_*L$ are locally constant, and the composition isomorphism identifies the latter with $(qp)_*L$. Pasting the two universally invertible base-change mates gives the mate for $qp$. Pullback stability follows by absorbing an initial base change into the universal quantifier. This verifies the precise repair discussed in I ms.~79--81 and 117--123 without a blanket replacement of the word ``equivalence''.
\end{proof}

\begin{claim}{A.20}{Refuted}{Properness does not imply local constancy of fibres}
The inclusion $j:\{1\}\to[1]$ is proper in Grothendieck's presheaf convention, but its comma transition over $0\to1$ is $\varnothing\to e$. Thus the assertion about proper maps at I ms.~102, companion row Ib-C003, is false already for spaces and set coefficients.
\end{claim}
\begin{proof}
The only arrow issuing from the image object $1$ is its identity. Hence $j$ is a Grothendieck opfibration. Its strict-fibre inclusions into the over-comma categories are identities: over $0$ both are empty and over $1$ both are terminal. The same is true after any base change because the inclusion becomes a cosieve. This is precisely the proper fibre formula. Nevertheless the two comma categories are empty and terminal, so their transition is not a weak homotopy equivalence, and for sets it changes the number of connected components. The category has two objects, one nonidentity arrow, and forced composition; the checker verifies these comma categories. Empty fibres are essential to the counterexample and cannot be discarded.
\end{proof}

\subsection{What the finite fibre theory supplies, and what Chapter XIII adds}

\begin{claim}{A.21}{Proved}{Finite products and the one-fibrant-leg pullback comparison}
Assume the hypotheses of A.11 and slice replacement equivalences as in A.12. In the localized fibrant slice over $S$, the product of $X\to S$ and $Y\to S$ is represented by $X\times_S Y\to S$. Arbitrary derived reindexing preserves these finite products. Moreover, if $X\to S$ is intrinsic and $Y\to\overline Y\to S$ is a replacement with $Y\to\overline Y\in W$ and $\overline Y\to S$ intrinsic, then $X\times_S Y\to X\times_S\overline Y$ lies in $W$. Thus a strict pullback with one intrinsic leg computes the same relative homotopy object as replacing the other leg first. These are the formal assertions at XII ms.~241--248.
\end{claim}
\begin{proof}
The fibre product of two intrinsic fibrations is again intrinsic by composition and base-change closure. A map between two such products induced by $u:X\to X'$ and $v:Y\to Y'$ in $W$ factors as
\[
X\times_SY\longrightarrow X'\times_SY
 \longrightarrow X'\times_SY'.
\]
The first arrow is a pullback of $u$ along the intrinsic object $Y\to S$, and the second a pullback of $v$ along $X'\to S$, so both are in $W$. The product functor and diagonal are therefore relative functors, and their ordinary adjunction descends by A.4. This proves the product universal property, not merely an abstract comparison of its objects. The terminal object of the slice is $1_S$, and its inclusion from the point is also a relative right adjoint, so the terminal-object assertion descends. Pullback commutes canonically with these products before localization, hence after localization. The last displayed map is a direct test of the definition of $X\to S\in\operatorname{Fib}_W$.
\end{proof}

\begin{claim}{A.22}{Proved}{The intrinsic factorization content of the weak Quillen axioms}
The conditions M2, M$'$, and M$''$ of XIII ms.~8 imply both intrinsic factorizations: every arrow factors as an intrinsic cofibration followed by a $W$-map, and as a $W$-map followed by an intrinsic fibration. They consequently imply the universal-equivalence equalities of XIII Corollary~1, ms.~10, and the fibrewise assertion of XIII Proposition~4, ms.~10--11. They do not identify the designated classes $TC,TF$ with $C\cap W,F\cap W$.
\end{claim}
\begin{proof}
M2 provides factorizations using $(C,TF)$ and $(TC,F)$. M$'$ puts $TC,TF$ in $W$, while M$''$ puts $C$ and $F$ in the two intrinsic classes. This gives the asserted factorizations with their exact order. A.11 gives $\operatorname{Fib}_W\cap W=W^u$ and the fibrewise assertion; reversing arrows gives the cofibration equality. The hypotheses only contain inclusions $TC\subseteq C\cap W$ and $TF\subseteq F\cap W$. The manuscript later gives separate names to equality on the left and on the right. Its term ``exact'' for a weak Quillen structure imposes $W=TF\circ TC$, whereas ``exact \`a gauche'' imposes $TC=C\cap W$. These are different formulas and none may be substituted for another by terminology alone.
\end{proof}

\begin{claim}{A.23}{Proved}{The variance assertions of XIII section 4 are lifting adjunctions}
Let $\mathcal M(I)=\operatorname{Fun}(I^{\mathrm{op}},\mathcal M)$, suppose the ordinary Kan extensions used below exist, and let the two closed lifting pairs be $(C(I),TF(I))$ and $(TC(I),F(I))$, with $C(I),TC(I)$ defined objectwise. For $u:I'\to I$ one has
\[
u^*C(I)\subseteq C(I'),\quad u^*TC(I)\subseteq TC(I'),
\quad u_*F(I')\subseteq F(I),\quad u_*TF(I')\subseteq TF(I),
\]
and
\[
\begin{split}
u^*F(I)\subseteq F(I')&\Longleftrightarrow u_!TC(I')\subseteq TC(I),\\
u^*TF(I)\subseteq TF(I')&\Longleftrightarrow u_!C(I')\subseteq C(I).
\end{split}
\]
If every component of $i\backslash I'$ has an initial object and the relevant coproducts exist, then all four containments in the last two equivalences hold. In particular injective fibrant diagrams are objectwise fibrant. For a walking-arrow diagram $X_0\leftarrow X_1$, injective fibrancy is equivalent to $X_0$ fibrant and $X_1\to X_0\in F$; this also implies that $X_1$ is fibrant.
\end{claim}
\begin{proof}
Restriction preserves objectwise classes. For any adjunction $L\dashv R$, a commutative square with left edge $La$ and right edge $b$ corresponds bijectively to a square with left edge $a$ and right edge $Rb$, and the bijection takes fillers to fillers. Thus $La\perp b$ if and only if $a\perp Rb$. Applying this to the two adjunctions and using that the lifting pairs are closed proves every displayed implication and converse. If the comma components have initial objects, the left Kan extension is the coproduct of evaluations at those initial objects, since the diagram being colimited has presheaf variance. A left orthogonal class is closed under these coproducts: lifting out of a coproduct is exactly one independent lifting problem on each summand. Thus $u_!$ preserves $C$ and $TC$ objectwise. Taking $I'=e$ proves the evaluation assertion.

For the walking arrow, lift first at $0$ against a trivial cofibration, using fibrancy of $X_0$, and then lift at $1$ against $X_1\to X_0$. This proves sufficiency. Conversely, test diagrams supported only at $0$ to obtain fibrancy of $X_0$, and test a square whose component at $0$ is an identity to obtain the lifting property of $X_1\to X_0$. This proves the exact arrow criterion. The argument is the actual content of XIII Propositions~1--3, ms.~45--50, and requires no claim that ordinary $u_*$ preserves all weak equivalences.
\end{proof}

\begin{claim}{A.24}{Refuted}{Ordinary Kan extensions do not automatically descend}
For $(\mathrm{Cat},W_\infty)$, ordinary limits of finite diagrams need not preserve objectwise weak equivalences. Thus an ordinary $u_*$ in XIII ms.~45--50 cannot be identified with a derived $u_*$ without a derivability or fibrancy argument.
\end{claim}
\begin{proof}
Consider the cospan diagram $e\xrightarrow{0}[1]\xleftarrow{1}e$ and the natural transformation to the constant cospan $e\to e\leftarrow e$, identity at its two outer vertices and the unique functor $[1]\to e$ at the middle vertex. Each component is in $W_\infty$, because $[1]$ has a terminal object. The strict pullbacks are respectively the empty category and the terminal category, and their comparison is $\varnothing\to e\notin W_\infty$. All objects, arrows, and the empty fibre product are checked by the finite program. This is a counterexample to raw preservation, not a counterexample to the fibrant computation or to the derived adjunction.
\end{proof}

\begin{claim}{A.25}{Proved}{The noetherian sober Alexandrov qualification adds no infinite posets}
If a poset with its Alexandrov topology is noetherian and sober, the poset is finite. Hence the apparently more general admissibility qualification in XIII ms.~60--61, requiring the Alexandrov space of $I^{\mathrm{op}}$ to be noetherian and sober, introduces no genuinely infinite poset beyond the finite case.
\end{claim}
\begin{proof}
We use the specialization order, so closed sets are lower sets. Noetherianity forbids an infinite strictly decreasing chain of principal closed sets. It also forbids an infinite antichain: for pairwise incomparable points $x_n$, the finite unions of their smallest open neighbourhoods form a strictly increasing sequence of open sets. An infinite strictly increasing chain of principal closed sets has a union $C$ that is closed, because arbitrary unions of lower sets are lower. This union is irreducible: if it were contained in $F\cup G$ for two closed sets without being contained in either, choose chain members containing a point outside $F$ and a point outside $G$; a later principal closed member contains both and is irreducible, a contradiction. Sobriety gives $C=\overline{\{x\}}$. The point $x$ belongs to some member of the chain, which then contains all of $C$, a contradiction to strict increase. Thus there is no infinite chain in either direction and no infinite antichain. Every infinite poset has a countably infinite subset containing an infinite chain or antichain, by the infinite Ramsey theorem applied to comparability; within an infinite chain, the usual monotone-subsequence argument gives an infinite increasing or decreasing subsequence. Therefore the poset is finite. Reversing the order does not change this conclusion.
\end{proof}

\begin{claim}{A.26}{Refuted}{The intrinsic finite floor does not provide all small Kan extensions}
Finite limits, finite colimits, two-out-of-three, and both intrinsic factorization axioms do not imply existence of all small homotopy Kan extensions. A counterexample is the category of finite sets with isomorphisms as $W$.
\end{claim}
\begin{proof}
Every map is both an intrinsic fibration and an intrinsic cofibration, since pullback and pushout preserve isomorphisms. Every map therefore has both required factorizations, inserting an identity as the $W$-factor. All finite limits and colimits exist. Localization changes nothing. But the diagonal functor $\mathrm{FinSet}\to\mathrm{FinSet}^{\mathbb N}$ has no right adjoint at the constant family of two-element sets: a representing finite set would have $2^{\aleph_0}$ global elements by the product universal property. It has no left adjoint at the constant singleton family either: a representing finite set $L$ would have $2^{\aleph_0}$ maps to a two-element set, whereas $\operatorname{Hom}(L,2)$ is finite. Thus XII's finite diagram theory and XIII's finite admissibility statement cannot be promoted to an all-small-shapes derivator without additional hypotheses.
\end{proof}

\begin{claim}{A.27}{Proved}{Closedness upgrades historical exactness to both trivial-class equalities}
Let $(C,TC,F,TF)$ satisfy the weak Quillen axioms of XIII ms.~27--28, including two-out-of-three for $W$, the two factorizations, and $TC\subseteq C\cap W$, $TF\subseteq F\cap W$. If both lifting pairs are closed and $W=TF\circ TC$, then
\[
TC=C\cap W,\qquad TF=F\cap W.
\]
Moreover $W$ is retract-closed. With the finite limits and colimits of M0 these data are a closed model structure in the original finite-limit convention. The two intrinsic containments M$''$ make this structure proper. Thus a \emph{closed exact} structure, as required by the uniform hypothesis at XIII ms.~57, is stronger than the unqualified exact structure distinguished in A.22.
\end{claim}
\begin{proof}
For $i:A\to B$ in $C\cap W$, exactness supplies $i=pj$ with $j:A\to Z$ in $TC$ and $p:Z\to B$ in $TF$. Since $i\perp p$, the square with top $j$ and bottom $1_B$ has a lift $s:B\to Z$. The identities $si=j$, $ps=1_B$ display $i$ as a retract of $j$ in the arrow category. A left orthogonal class is retract-closed, so $i\in TC$. The dual argument for $p\in F\cap W$ proves the second equality.

For completeness we give the retract argument for $W$, which does not require functorial factorizations; compare \citet[Lemma~11.3.3]{BISITERiehl2014}. First, if $f\in F$ is a retract of $w\in W$, factor $w=vj$ with $j\in TC$, $v\in TF$. Lifting $j$ against $f$ extends the retraction of their sources and displays $f$ as a retract of $v$, hence $f\in TF$. For a general retract $f:A\to B$ of $w:X\to Y$, write the sections as $a:A\to X$, $b:B\to Y$ and retractions as $r:X\to A$, $s:Y\to B$. Factor $f=hg$, with $g:A\to Z$ in $TC$ and $h:Z\to B$ in $F$. Form $P=X\amalg_A Z$, with structure maps $i:X\to P$ and $x:Z\to P$. Then $i\in TC$. The maps $w,bh$ induce $j:P\to Y$ with $ji=w$, so $j\in W$. The maps $gr,1_Z$ induce $y:P\to Z$ with $yx=1_Z$ and $hy=sj$. Hence $(x,b)$ and $(y,s)$ display $h$ as a retract of $j$. The first case gives $h\in W$; therefore $f=hg\in W$. The remaining model axioms are precisely the stipulated lifting, factorization, and finite-limit axioms. Finally M$''$ says directly that pushout along $C$ and pullback along $F$ preserve $W$, which is properness.
\end{proof}

\begin{claim}{A.28}{Proved}{The uniform theorem of XIII and the exact admissibility hypothesis}
Assume exactly the uniform-left hypothesis of XIII ms.~57: for every $I$ in the specified diagram domain, the objectwise classes $C(I),TC(I)$ and their right orthogonals $TF(I),F(I)$ form a closed exact weak Quillen structure for the objectwise class $W(I)$. Let $u:I'\to I$ be a functor in the domain for which the ordinary presheaf right Kan extension $u_*$ exists. Then $u_*$ preserves weak equivalences between fibrant diagrams. It consequently induces a right adjoint
\[
u^*: \mathcal M(I)[W(I)^{-1}] \rightleftarrows\
\mathcal M(I')[W(I')^{-1}]:\mathbf R u_*.
\]
If the comma categories $I'/i$ and their inclusions into $I'$ belong to the diagram domain, the canonical pointwise formula is
\[
i^*\mathbf R u_*X\ \cong\ \mathbf R\!\lim_{(I'/i)^{\mathrm{op}}}X|_{I'/i}.
\]
In particular every functor between finite posets is admissible under the source's hypotheses. No restriction to finite posets is necessary for admissibility once the relevant ordinary Kan extension and uniform diagram structures exist. In the historical numbering, this proves the right-hand Der~4 and Der~5 assertions of XIII ms.~61; these numbers do not mean modern Der~4 and strongness Der~5.
\end{claim}
\begin{proof}
A.27 applies separately on every diagram category. A.23 shows that $u^*$ preserves $C$ and $TC$ and $u_*$ preserves $F$ and $TF$. On the full subcategory of fibrant diagrams, terminal maps and products are fibrant, and factorization of an arrow as $TC$ followed by $F$ stays inside that subcategory. Apply the elementary functor lemma proved in A.17 to $u_*$: its hypothesis holds because $F\cap W=TF$ and $u_*$ preserves $TF$. Thus $u_*$ preserves every weak equivalence between fibrant objects. This is Brown's lemma, not an additional infinite-product or transfinite-closure assumption.

We recall explicitly the homotopy-category calculation used to derive the adjunction. For a closed model structure with finite limits and colimits, a cofibrant source $A$ and fibrant target $B$ satisfy
\[
\operatorname{Hom}_{\mathcal M[W^{-1}]}(A,B)
\cong \operatorname{Hom}_{\mathcal M}(A,B)/\!\sim,
\]
where $\sim$ is the common left- and right-homotopy relation. Cofibrant and fibrant replacements induce equivalences from their full subcategories after localization. These are the original finite-limit form of Quillen's homotopy-category theorem, I.1, Theorem~1; no functorial replacement is needed \citep{BISITEQuillen1967}. To apply its derived-adjunction argument, take $A$ cofibrant in $\mathcal M(I)$ and $B$ fibrant in $\mathcal M(I')$. The ordinary adjunction identifies maps $u^*A\to B$ with maps $A\to u_*B$. It identifies their homotopy classes: a cylinder of $A$ is taken to a cylinder of $u^*A$, since restriction preserves cofibrations and all objectwise weak equivalences; homotopies out of that cylinder correspond under the original adjunction. The homotopy-category formula therefore gives a natural bijection on localized Hom sets. Replacing arbitrary sources and targets by cofibrant and fibrant ones gives the asserted adjunction. Preservation on fibrant objects, already proved, makes the computation by $u_*B$ independent of the chosen fibrant representative. This is Quillen I.4, Theorem~3, with the comparison maps specified.

For the pointwise assertion, write $q_i:I'/i\to I'$. Its restriction preserves fibrant diagrams by A.23: for $j\in I'$, every component of $j\backslash(I'/i)$ has an initial object, namely the factorization beginning with $1_j$ of a fixed arrow $u(j)\to i$. Thus a fibrant $B$ restricts to a fibrant diagram on $I'/i$. The ordinary formula $(u_*B)(i)=\lim_{(I'/i)^{\mathrm{op}}}q_i^*B$ computes both derived sides, and it computes their canonical mate. Finite posets and their relevant comma categories require only finite limits, supplied by M0. A.25 explains why the source's noetherian sober qualification adds no infinite case.
\end{proof}

\subsection{Formal closure corrections and the absolute localizer axioms}

\begin{claim}{A.29}{Proved}{The polarity calculation in II section 5 requires closure after a join}
For any relation $R\subseteq P\times Q$, the maps
$A\mapsto A^\perp=\{q\mid pRq\text{ for all }p\in A\}$ and its transpose define an antitone Galois connection on the power sets. Double orthogonality is an extensive idempotent closure; closed subsets are closed under arbitrary intersections, and their join is the closure of their union. If $P,Q$ are complete posets and the principal-polar condition (D$_0$) of II ms.~69 holds, write $c$ for the induced closure on $P$. For a subset $A\subseteq P$, its polar closure is the principal lower set
\[
A^{\perp\perp}=\mathord\downarrow c\!\left(\bigvee A\right)
=\mathord\downarrow c\!\left(\bigvee_{a\in A}c(a)\right).
\]
The final closure cannot be omitted from II formula~(5.20), ms.~76. One must also distinguish the element $c(\bigvee A)$ of $P$ from its principal lower set, a subset of $P$.
\end{claim}
\begin{proof}
The equivalence $B\subseteq A^\perp\Longleftrightarrow A\subseteq B^\perp$ is precisely the assertion that every pair in $A\times B$ belongs to $R$. Extensivity and antitonicity give the threefold identities and hence idempotence, as in A.6. An intersection of closed subsets is closed because it is the polar of the union of their opposite polars. The least closed set containing a union is its closure. Under (D$_0$), taking all $a\in A$ in the relation amounts to taking their supremum; its polar and double polar are the stated principal lower sets. Extensivity gives $\bigvee A\leq\bigvee c(A)$, while monotonicity gives $c(a)\leq c(\bigvee A)$ for every $a$. Applying $c$ to the two inequalities proves the second equality.

For the omitted-closure issue, take the usual annihilator polarity between subsets of $V=\mathbb F_2^2$ and subsets of its dual, ordered by inclusion. The induced closure on subsets of $V$ is linear span. The two one-dimensional subspaces $\{0,(1,0)\}$ and $\{0,(0,1)\}$ are closed. Their union has three elements and is not closed; the correct closed join is all four elements of $V$. The finite checker evaluates all scalar products and double annihilators. This verifies the correction without relying on the companion's wording.
\end{proof}

\begin{claim}{A.30}{Refuted}{Arbitrary unions do not preserve composition closure}
The class of arrow classes containing identities and closed under both pullback and composition is not closed under arbitrary unions, contrary to the preliminary lattice description in II ms.~65--67. Intersections have that property; a join requires taking the closure of the union under the specified operations.
\end{claim}
\begin{proof}
In the chain $0<1<2$, let $L_1$ consist of the identities and $0\to1$, and let $L_2$ consist of the identities and $1\to2$. Each is pullback-closed: pulling the nonidentity arrow back along an arrow to its target gives either itself or an identity. Each composes, since no two nonidentity arrows in one class are composable. Their union contains $0\to1$ and $1\to2$ but not $0\to2$. The complete category and the closure tests are checked in the attached finite program. More generally, intersections preserve each closure condition because every required output lies in every intersectand; the smallest closed class containing a union is therefore the intersection of all closed classes containing it.
\end{proof}

\begin{claim}{A.31}{Proved}{The bi-site remark fixes a question rather than a general theorem}
At VI ms.~81--84 the data $(\mathcal M,W)$ specify the coarse slice localization, but the proposed fine classes require additional test data. The two comma tests in $\mathrm{Cat}$ have opposite variances: $X/s$ varies covariantly with $s$, and $s\backslash X$ contravariantly. In XII ms.~222 the class labelled $W^g$ uses the former tests and its universal preservation class is the smooth class discussed in Part B; XII ms.~224's $W^d$ uses the latter. The letters on an earlier page are not a substitute for this variance check.

The remark that weak globular infinity-groupoids might have no useful localization theory is a suggestion at VI ms.~83--84. It is not a nonexistence theorem about fine relative theories. The attached 2026 claims about higher categories and globular groupoids remain ``claimed (OpenAI, 2026, unrefereed)'' here \citep{OpenAI2026Thomason,OpenAI2026Groupoids}; no conclusion of Part A depends on either.
\end{claim}
\begin{proof}
An arrow $s\to t$ postcomposes $X/s\to X/t$ and precomposes $t\backslash X\to s\backslash X$. The slice projection $S/s\to S$ is a local isomorphism in the discrete-fibration sense and its pullback is $X/s$; the opposite projection gives the other test. XII ms.~222 and 224 explicitly use these respective classes. The equality with the relevant smooth or proper preservation class needs the extra categorical argument of Part B and does not follow merely from renaming a Galois transporter. VI ms.~83 describes the two slice operations and the involution, and uses speculative language for the groupoid remark; it supplies no axiom forbidding a chosen finer class. Ordinary groupoids provide the unconditional test case in Part E.
\end{proof}

\begin{claim}{A.32}{Refuted}{The literal absolute package does not imply relative Theorem A}
The literal axioms Loc~(1,2,3) of XVI, with the moderate saturation of A.1, do not imply Loc~(3 bis). The failure already occurs in a triangle of finite groupoids over the one-object category $BC_2$. Thus the absolute package left undecided in POLES A.16 must be distinguished from the relative package used throughout the present paper.
\end{claim}
\begin{proof}
For a functor $f:A\to B$ between arbitrary small categories define
\[
 f\in W^\dagger\quad\Longleftrightarrow\quad
 \pi_0(f)\text{ is bijective, and }
 \bigl(\lvert\pi_0(A)\rvert=1\Longrightarrow f\in W_\infty\bigr).
\]
Thus on connected categories we retain exactly the weak homotopy equivalences; on categories with at least two components we require only a bijection of components. The empty-to-empty functor also belongs to $W^\dagger$.

We first verify the axioms on all small categories, including infinitely many components. Isomorphisms belong to $W^\dagger$. If two of $f,g,gf$ belong to $W^\dagger$, two-out-of-three for bijections shows that all three maps induce bijections on components, so their three categories have the same component cardinality. If it is not one, no additional condition is required; if it is one, two-out-of-three for $W_\infty$ gives the assertion. For moderate saturation, suppose $f:A\to B$ and $g:B\to A$ have $gf,fg\in W^\dagger$. Bijections of both composites imply that $\pi_0(f)$ and $\pi_0(g)$ are bijections. Again only the connected case requires more. In that case both composites are weak homotopy equivalences, and hence both factors are weak homotopy equivalences: their images in the homotopy category of spaces have invertible composites in both orders, which gives a left and a right inverse for each factor. Classifying spaces are CW complexes, so isomorphism there is equivalent to weak homotopy equivalence. This proves precisely Loc~(1), without assuming full arrow-retract closure.

A category is $W^\dagger$-aspherical if and only if its nerve is weakly contractible. Indeed its map to $e$ must have exactly one component, and the remaining condition is exactly membership in $W_\infty$. A category with a terminal object has contractible nerve, proving Loc~(2). If every $A/b$ is $W^\dagger$-aspherical, each has weakly contractible classifying space; these CW complexes are contractible. Quillen's Theorem~A, in its dual comma formulation, gives $f\in W_\infty\subseteq W^\dagger$, proving the absolute Loc~(3) \citep[Theorem~A, p.~93]{BISITEQuillen1973}.

For the failure of the relative clause, use additive notation $G=\mathbb Z/2\mathbb Z$. Let $C=BG$ have one object $c$, morphism group $G$, identity $0$, and composition addition. Let $E$ be the groupoid with objects $x\in G$ and exactly one arrow
\[
 (x,g):x\longrightarrow x+g\qquad(g\in G).
\]
Its identity at $x$ is $(x,0)$ and
$(x+g,h)\circ(x,g)=(x,g+h)$. Define $p:E\to C$ by $p(x)=c$, $p(x,g)=g$. Let $K=BG$, let $A=E\times K$, let $B=E$, and let $u:A\to B$ be projection, with structure maps $pu:A\to C$ and $p:B\to C$. Explicitly, $A$ has the same two object labels as $E$, an arrow $(x,g,k):x\to x+g$ for each $g,k\in G$, identities $(x,0,0)$ and composition
\[
 (x+g,h,l)\circ(x,g,k)=(x,g+h,k+l).
\]
The functor $u$ forgets $k$. This gives all objects, arrows, identities, compositions and functors in the triangle.

The groupoid $E$ is connected and has exactly one arrow between any two objects; it is equivalent to $e$. The groupoid $A$ is connected with automorphism group $G$ at either object, and is equivalent to $BG$. Consequently $u$ induces $G\to1$ on fundamental groups. It is not in $W_\infty$, and since both categories are connected it is not in $W^\dagger$.

There is just one object $c$ of the base, so the relative premise has one comma comparison. The objects of $E/c$ are pairs $(x,t)\in G^2$, where $t:p(x)\to c$. An arrow labelled $g$ from $(x,t)$ to $(x+g,t')$ exists precisely when $t'+g=t$. Thus $t+x$ is invariant under arrows. For each of its two values there are two objects and exactly one arrow between every pair. Hence $E/c$ has exactly two connected components, each equivalent to $e$. There is a strict identification
\[
 A/c=(E/c)\times K,
\]
under which $u/c$ is projection. Its source also has exactly two components, and it is a bijection on those components, indexed on both sides by $t+x$. Therefore $u/c\in W^\dagger$, whereas $u\notin W^\dagger$. This refutes Loc~(3 bis).

The saturation boundary is explicit. The nonweak map $BG\to e$ is a retract, in the arrow category, of
\[
 (BG\amalg e)\longrightarrow(e\amalg e),
\]
which collapses the first component and is the identity on the second. This latter arrow is in $W^\dagger$. The sections are the inclusions of the first components; the source retraction is the identity on $BG$ and sends the added point to its unique object, and the target retraction is the unique map to $e$. These maps commute and their retraction composites are identities. Thus $W^\dagger$ fails Loc~(1 bis); the example proves the asserted failure for the actual literal Loc~(1), not for an unmentioned stronger saturation convention. The modern split-retraction clause nevertheless follows from the already verified moderate condition by taking $ri=1$.

The program \texttt{code/a\_absolute\_relative.py} independently verifies all finite category laws, the functors, both comma categories, their components and automorphism groups, and the displayed arrow retract. The general axiom check above does not rest on any finite sampling. For comparison, a base with a terminal object cannot give this kind of counterexample: its terminal-object comma comparison is already the original functor. In particular $[1]$ is not an unresolved test of the relative implication.
\end{proof}

\begin{claim}{A.33}{Refuted}{The converse cocartesian assertion of XII ms.274 needs conservativity}
Theorem~3$^\circ$ of XII ms.~274 is false without the conservativity premise introduced in its part~2$^\circ$. The assertion that for a map of homotopy cocartesian squares $f_0,f_3\in W$ imply $f_1,f_2\in W$ fails even when $W$ is strongly saturated and the intrinsic cofibration factorization exists. Corollary~1 on that page is also false without that extra premise.
\end{claim}
\begin{proof}
Take $\mathcal M=\mathrm{FinSet}$ and $W$ the isomorphisms. All maps are intrinsic cofibrations, every map factors as itself followed by an identity, and the $W$-cocartesian squares are the ordinary pushout squares. Let $D=\{0,1\}$ and $e=\{*\}$. The two spans
\[
 X=(D\xrightarrow{1_D}D,\ D\longrightarrow e),\qquad
 Y=(D\longrightarrow e,\ D\longrightarrow e)
\]
both have pushout $e$. The maps $f_0=1_D$, $f_1:D\to e$, and $f_2=1_e$ define a map of spans and induce $f_3=1_e$. Thus $f_0,f_3\in W$ while $f_1\notin W$. All four component maps and the pushout equivalence relations are checked by the finite program.

Here a $W$-aspherical set is a singleton. In any pushout square whose initial and final vertices are singletons, both other vertices are singletons: the maps from the initial singleton make both nonempty, and identifying just their chosen basepoints cannot collapse any additional element. Thus both conditions~1 and~1$'$ of Corollary~1 hold in this example, although its claimed equivalent implication fails. The counterexample therefore locates the missing hypothesis rather than merely a gap in its proof. The source's Corollary~2, ms.~275, explicitly restores conservativity. In the preceding Cons~2, ms.~273, ``cart\'esien'' must read ``cocart\'esien'' for the stated reduction; this word is printed that way in the edition, not introduced by transcription.
\end{proof}

\begin{claim}{A.34}{Refuted}{Stabilizers alone do not define the antitone correspondence of the cancelled section}
For arbitrary arrow classes $A$, the operation $A\mapsto\operatorname{Stab}(A)$ is not antitone. Thus the first, cancelled proposal in II section~4 cannot be used in place of the corrected transporter construction in section~6.
\end{claim}
\begin{proof}
In $[2]$, let $A$ consist of the identities and $0\to2$, and let $A'$ be all arrows. Both contain identities and compose, and $A\subseteq A'$. The arrow $1\to2$ does not stabilize $A$, because the pullback of $0\to2$ is $0\to1\notin A$. It does stabilize $A'$, as every arrow does. Therefore $\operatorname{Stab}(A')\nsubseteq\operatorname{Stab}(A)$. The checker verifies both stabilizers by all slice tests. By contrast, $T_W(A)$ with fixed target class $W$ is antitone because enlarging $A$ only adds tests; that is the correction proved in A.6.
\end{proof}
